\chapter{Related Work}
\label{ch:related-work}

This chapter reviews four bodies of work that this thesis draws on: preference-based
alignment and how preference data becomes a training signal
(Section~\ref{sec:rw-rlhf-rlaif}); the reliability of AI judges when they replace
human annotators (Section~\ref{sec:rw-judge-reliability}); the granularity of
feedback in long-form generation, which is the variable this thesis manipulates
(Section~\ref{sec:rw-finegrained}); and automatic factual-consistency evaluation,
which supplies the judge (Section~\ref{sec:rw-factual-eval}).
Section~\ref{sec:rw-synthesis} draws these together into the specific gap this
thesis addresses, and Section~\ref{sec:rw-positioning} states what separates the
present work from the studies closest to it.

\section{Preference-Based Alignment: RLHF, RLAIF, and Offline Optimization}
\label{sec:rw-rlhf-rlaif}

Preference-based alignment collects pairwise comparisons between candidate model
outputs and converts those comparisons into a training signal. The approach
originates in preference learning for reinforcement learning, where human
annotators compare short trajectory segments and a reward model is fitted to
reproduce their comparisons, allowing an agent to be trained on objectives that
are easy to judge but hard to specify numerically \citep{christiano2017deep}. The
choice to elicit comparisons rather than absolute scores is deliberate: relative
judgments between two concrete outputs are more consistent between annotators than
absolute ratings on an abstract scale, and the Bradley--Terry model
(Chapter~\ref{ch:background}, Section~\ref{sec:bg-bradley-terry}) provides a
principled way to recover a scalar reward from them.

\citet{stiennon2020learning} adapted this recipe to summarization, and their
pipeline is the direct ancestor of the one used in this thesis. They collect human
comparisons between candidate summaries, train a reward model on those comparisons,
and then fine-tune a policy against the learned reward with reinforcement learning.
Two of their findings matter here. First, the reward model's own held-out agreement
with human preferences is reported as a diagnostic in its own right, which
establishes the precedent for treating reward-model accuracy as a meaningful
measurement rather than merely an intermediate artifact, as this thesis does.
Second, they observe that optimizing too hard against a learned reward degrades
true quality, an early instance of the reward-overoptimization problem that
motivates careful attention to the quality of the preference signal itself. RLHF
was subsequently scaled to general instruction-following \citep{ouyang2022training},
establishing the reward-model-plus-policy-optimization pattern as the default
alignment recipe.

Because human comparison data is expensive, reinforcement learning from AI feedback
(RLAIF) substitutes a model for the human annotator. Constitutional AI generates
critiques and revisions from a written set of principles and derives preference
labels from them, removing human annotation from the harmlessness training loop
almost entirely \citep{bai2022constitutional}. \citet{lee2023rlaif} compare RLAIF
and RLHF directly on summarization and dialogue tasks and report broadly
comparable downstream performance, which is the empirical result that makes
AI-generated preference data a credible substitute rather than a compromise. They
also emphasize that the outcome depends substantially on how the judge is
prompted and calibrated, a dependency that recurs throughout
Section~\ref{sec:rw-judge-reliability} and, in a different form, in the
evaluator-configuration effects observed here
(Chapter~\ref{ch:results}, Section~\ref{sec:results-mode-sweep}).

On the optimization side, RLHF pipelines have typically used on-policy algorithms
such as Proximal Policy Optimization \citep{schulman2017ppo}, which require
sampling from the policy during training and maintaining a separate reward model.
Direct Preference Optimization (DPO) removes both requirements by reparameterizing
the RLHF objective so that the optimal policy can be recovered in closed form
directly from preference pairs \citep{rafailov2023dpo}. DPO was the objective
originally proposed for this thesis, chosen because it isolates the effect of the
preference data from the dynamics of reinforcement learning. As
Chapter~\ref{ch:methodology} describes, DPO was implemented and produced early
results before being dropped from the final comparison in favor of evaluating
Bradley--Terry reward models directly. That change narrows the comparison further
than DPO alone would: it removes both the optimization dynamics that motivated
adopting DPO and the decoding variance introduced by generating and scoring a
fine-tuned policy's outputs, leaving a more targeted question about whether the
preference data itself is more learnable.

\section{Reliability Risks in AI Supervision and LLM-as-a-Judge}
\label{sec:rw-judge-reliability}

Substituting a language model for a human annotator introduces a distinct class of
error. Benchmark studies of LLM-as-a-judge document sensitivity to prompt wording,
position bias favoring whichever candidate appears first, verbosity bias favoring
longer answers regardless of quality, inconsistent scoring across repeated calls,
and self-enhancement effects in which a judge prefers outputs from models similar
to itself \citep{zheng2023judging}. \citet{liu2023geval} document related
calibration problems in LLM-based evaluation of generated text, and a recent
survey synthesizes the accumulated evidence on these failure modes and the
mitigations proposed for them, including rubric enforcement, pairwise rather than
absolute evaluation, and position swapping \citep{gu2024survey}.

These issues are more consequential in RLAIF than in evaluation. An unreliable
evaluation metric produces a noisy estimate of a model's quality; an unreliable
judge used to construct training data produces systematically mislabeled
supervision, and the model is then optimized to fit that mislabeling.
\citet{casper2023open} catalog failure modes of this kind across the RLHF
pipeline, including reward misspecification and the divergence between training
loss and intended behavior, and argue that problems in the preference data
propagate into the trained policy in ways that are difficult to detect after the
fact.

One response to this literature is to treat the judge as a fallible instrument and
wrap it in a reliability protocol: confidence gating, order randomization,
evidence-grounded rationale requirements, and prompt-robustness checks. The design
adopted here removes the exposure at its source instead, by using a fixed,
deterministic factuality model rather than a generative judge
(Section~\ref{sec:rw-factual-eval}). A deterministic model cannot exhibit prompt
sensitivity, position bias, or run-to-run inconsistency, because it is never
prompted and returns no natural-language output. This eliminates the category of
risk the protocol above is designed to manage, at two costs: it does not make the
evaluator correct, and it removes the rationale text that a human audit would
examine. Chapter~\ref{ch:methodology}, Section~\ref{sec:methodology-scope} states
what this implies for the claims the thesis can support.

\section{Feedback Granularity and Credit Assignment in Long-Form Generation}
\label{sec:rw-finegrained}

The variable this thesis manipulates is how localized the supervision signal is.
A single preference over a multi-sentence output provides one bit of information
about a text that may contain a dozen independently checkable claims, and it does
not indicate which claim caused the preference. This is a credit assignment
problem, and it is present even in the original preference-learning formulation:
\citet{christiano2017deep} have annotators compare short trajectory segments
rather than whole episodes precisely to make the resulting supervision denser and
easier to attribute.

The closest prior work to this thesis is \citet{wu2023finegrained}, who study
fine-grained human feedback for long-form question answering. Their approach
differs from holistic RLHF along two dimensions simultaneously. First, feedback is
collected at the density of individual sentences or sub-sentence spans rather than
whole responses, so an annotator marks the specific span that is irrelevant,
untruthful, or incomplete. Second, they train several distinct reward models, one
per error category, and combine their outputs into a dense reward signal during
reinforcement learning, so the policy receives a reward at each segment rather than
a single scalar at the end of the sequence. They report that this improves both the
resulting policy and the reward models themselves relative to holistic feedback,
and that the multiple reward models allow the relative weight of different
objectives to be adjusted after training.

Three differences separate their setting from this thesis. Their fine-grained
labels come from trained human annotators, whereas this thesis asks whether the
same localization strategy is still useful when the localized judgments come from
an automatic judge, which is cheaper by orders of magnitude but noisier per label.
Their fine-grained signal is consumed as a dense per-segment reward during
reinforcement learning, whereas GCA aggregates local scores back into a single
summary-level preference pair so that the resulting dataset remains compatible
with standard pairwise preference optimization, requiring no change to the training
objective. And they vary feedback density and objective decomposition together,
whereas this thesis holds the objective fixed at factual consistency alone and
varies only granularity, which isolates the granularity question at the cost of
addressing a narrower one.

Three more recent lines of work place sentence- or step-level signals at
other points in the alignment pipeline, and each should be distinguished
carefully from the present study, since none leaves the question examined
here open in general.

\citet{gu2025maskdpo} introduce Mask-DPO, which incorporates sentence-level
factuality signals directly into the DPO objective. Rather than treating a whole
response as preferred or dispreferred, Mask-DPO uses sentence-level factuality
information to learn only from the factually correct sentences within a preferred
response, and to avoid penalizing factually correct content that appears inside a
dispreferred one. The motivation is that response-level factuality supervision is
noisy, because preferred and dispreferred responses each mix accurate and
inaccurate sentences. The intervention point is the policy-optimization objective:
the preference pair remains a response-level object, and the granular signal is applied as a mask during training. The present thesis intervenes earlier and
leaves the training objective untouched. It asks how automatic factuality
judgments at different granularities should be converted into the pairwise label
in the first place, and measures the learnability of the resulting labels under a
fixed Bradley--Terry architecture.

\citet{qiu2025sentencelevel} modify the reward model itself, proposing a
sentence-level reward model that assigns scores at sentence boundaries and
aggregates them into a response-level score, while still learning from
response-level human preferences. Their aim is to address reward sparsity by
producing an intermediate granularity between response-level and token-level
reward modeling. The distinction from this thesis is architectural: they change
what the reward model computes, and the preference labels they learn from are
conventional response-level human judgments. This thesis does not modify the
reward-model architecture at all. The Bradley--Terry model, its backbone, and its
training procedure are held fixed across both conditions precisely so that
differences in outcome can be attributed to the preference labels rather than to
the model consuming them.

Closer still to GCA's decompose-then-aggregate structure is process reward
modeling for multi-step reasoning \citep{lightman2023verify}. Rather than
training a single model to judge a complete solution as correct or
incorrect, \citet{lightman2023verify} have human annotators label the
correctness of each individual reasoning step in a chain-of-thought
solution, train a reward model to reproduce those per-step judgments, and
show that this process-level supervision outperforms a response-level
outcome reward for selecting correct solutions among many sampled
candidates. The structural resemblance to GCA is real: both decompose a
long output into smaller units, score each unit, and use the per-unit
scores rather than a single holistic judgment. Two differences remain.
Process reward models keep the per-step scores as the object consumed
downstream, typically for reranking or search over sampled solutions,
whereas GCA aggregates its per-sentence scores back into one summary-level
score so that the output is a conventional pairwise preference, compatible
with unmodified Bradley--Terry training. And the per-step labels in
\citet{lightman2023verify} come from trained human annotators judging
mathematical correctness, a domain with a comparatively well-defined notion
of a correct step, whereas GCA's per-sentence labels come from an automatic
factuality judge applied to open-ended news summaries, where what counts as
a faithful sentence is judged by a model rather than verified against a
solution. The result that step-level supervision outperforms outcome-level
supervision in their setting is consistent with the motivation for GCA, but
it is evidence from a different domain, a different label source, and a
different point of consumption in the pipeline, not a direct precedent for
the automatic, aggregated case this thesis studies.

Reward models are also increasingly evaluated in their own right rather than only
as components of a larger pipeline. RewardBench assembles prompt-chosen-rejected
trios across several domains and evaluates reward models by how often they rank
the preferred response above the dispreferred one
\citep{lambert2024rewardbench}. That pairwise-ranking accuracy is the same
quantity used as the dependent variable in this thesis
(Chapter~\ref{ch:methodology}, Section~\ref{sec:methodology-scope}), which places
the measurement on established ground; the difference is that RewardBench varies
the reward model against fixed data, whereas this thesis varies the data
construction against a fixed reward model.

\section{Factual Consistency in Summarization and Evaluation}
\label{sec:rw-factual-eval}

News summarization remains a standard benchmark domain, with CNN/DailyMail widely
used for both training and evaluation \citep{hermann2015teaching,
nallapati2016abstractive}. Abstractive systems have long produced fluent output
that nonetheless contradicts or embellishes the source \citep{see2017get}, and this
persistent gap between fluency and faithfulness is what makes factual consistency
a worthwhile optimization target rather than a solved preprocessing concern.

Automatic factuality evaluation has developed along two main lines. Entailment-based
methods treat the source as a premise and the summary as a hypothesis: FactCC
trains a classifier on synthetically corrupted summaries to detect factual
conflicts \citep{kryscinski2020evaluating}. Question-answering-based methods
instead generate questions from the summary and check whether the source yields
the same answers, an approach developed in FEQA \citep{durmus2020feqa} and refined
in QAFactEval, which systematically compares components of the QA-based pipeline
\citep{fabbri2022qafacteval}.

\paragraph{Decomposition and aggregation in SummaC.} The work most directly
relevant to GCA's design is SummaC \citep{laban2022summac}, which observes that
natural language inference models trained on sentence-level data transfer poorly
to document-level consistency checking, partly because the inputs are far longer
than anything seen in training. Their solution is to decompose the problem: they
compute an entailment matrix between each source sentence and each summary
sentence, then aggregate that matrix into a single score. Critically for this
thesis, they show that \emph{how} the matrix is aggregated changes performance
substantially. Their zero-shot variant takes, for each summary sentence, the
maximum entailment score across all source sentences, and then averages these
maxima across summary sentences. Their learned variant instead trains a
convolutional layer over the distribution of entailment scores per summary
sentence, retaining information that the maximum discards, and performs better.

That aggregation strategy is a design decision with measurable consequences is
therefore established in the evaluation literature before this thesis takes it up.
The contribution here is to ask the same question one step upstream, at the point
where scores are converted into training labels rather than into an evaluation
metric, and Chapter~\ref{ch:results}, Section~\ref{sec:results-formula-opt} finds
the same lesson holds: the choice between a mean and a penalty-weighted
aggregation changes the outcome more than the decision to decompose at all.

\paragraph{AlignScore.} The judge used throughout this thesis, AlignScore, was not
part of the literature surveyed in the original proposal. \citet{zha2023alignscore}
recast a broad collection of natural language processing tasks, including natural
language inference, question answering, paraphrase detection, and fact
verification, into a single ``information alignment'' format, and train one model
on the union. The resulting function takes a context and a claim and returns a
scalar in $[0,1]$ estimating whether the context supports the claim. Because it is
trained across many task formats rather than on entailment data alone, it
generalizes across factuality benchmarks better than the single-task methods above.

AlignScore also inherits the decomposition question from SummaC, and resolves it
internally: in its split evaluation modes, the context is divided into chunks and
the claim into sentences, and the final score aggregates the resulting matrix by
taking a maximum over chunks and then a mean over claim sentences, which is
structurally the same rule as SummaC's zero-shot variant. This internal behavior
turns out to matter a great deal for the present work.
Chapter~\ref{ch:results}, Section~\ref{sec:results-mode-sweep} finds that whether
GCA outperforms holistic scoring depends on which AlignScore mode is used, and
Chapter~\ref{ch:discussion}, Section~\ref{sec:disc-interpretation} argues that this
is because the split modes already perform, inside the judge, much of the
decomposition that GCA performs outside it.

\paragraph{Fine-grained evaluation.} More recent work evaluates summaries at
sub-summary granularity rather than assigning a single score. FineSurE assesses
faithfulness at the sentence level and content coverage against a list of extracted
key facts, producing a diagnostic profile rather than one number
\citep{song2024finesure}. Like SummaC, and unlike this thesis, it is positioned as
an evaluation framework: it tells a practitioner where a summary went wrong, but
does not convert that localized judgment into supervision for training.

FActScore decomposes a generation into individual atomic facts and checks
each against a knowledge source, then reports the fraction supported as a
single precision score \citep{min2023factscore}. Its decompose-score-aggregate
structure is the closest evaluation-side analogue to GCA's own aggregation
step (Chapter~\ref{ch:methodology}, Section~\ref{sec:bg-gca}): both split a
generation into small checkable units, score each independently, and combine
the results into one number. The differences are the same ones that
separate this thesis from SummaC and FineSurE. FActScore decomposes into
atomic facts extracted from the generation itself and verifies them against
an external knowledge source, whereas GCA decomposes into sentences and
scores each against the source article with a single fixed judge; and
FActScore's aggregate score is reported as a diagnostic metric, not
converted into a pairwise preference for training a reward model. The
aggregation choice it makes, a simple fraction of atomic facts supported,
is nonetheless a further data point for the same lesson this thesis draws
in Section~\ref{sec:results-formula-opt}: how per-unit scores are combined
is a design decision with measurable consequences, examined here independently
by two unrelated methods on two different granularities of decomposition.

\section{Synthesis: The Gap This Thesis Addresses}
\label{sec:rw-synthesis}

Table~\ref{tab:rw-gap} summarizes the closest lines of work along the dimension
that distinguishes the present study: the point in the pipeline at which
sentence-level information is introduced, and what each approach holds fixed while
doing so.

\begin{table}[htbp]
\centering
\caption[Where sentence-level information enters the pipeline in related work]{Where sentence-level information enters the pipeline in related work.}
\label{tab:rw-gap}
\begin{tabularx}{\textwidth}{lXX}
\toprule
Work & Sentence-level signal enters at & Held fixed \\
\midrule
\citet{stiennon2020learning} & not used & -- \\
\citet{lee2023rlaif}         & not used & -- \\
\citet{wu2023finegrained}    & reward signal during RL (human labels) &
  response-level objective \\
\citet{gu2025maskdpo}        & the DPO objective, as a sentence mask &
  preference pair construction \\
\citet{qiu2025sentencelevel} & the reward-model architecture &
  preference labels (human, response-level) \\
\citet{lightman2023verify}   & the reward model's own output (per-step) &
  training objective; labels (human, per-step) \\
\citet{laban2022summac}      & evaluation metric only & -- \\
\citet{song2024finesure}     & evaluation metric only & -- \\
\citet{min2023factscore}     & evaluation metric only & -- \\
\citet{zha2023alignscore}    & internal to the evaluator & -- \\
\midrule
This thesis & \textbf{construction of the preference label} &
  reward-model architecture and training procedure \\
\bottomrule
\end{tabularx}
\end{table}

Table~\ref{tab:rw-gap} organizes this literature by \emph{where} sentence-level
information is injected, which is the dimension along which the present study is
distinguished. Sentence-level factuality signals are clearly not unexplored: they
have been used to mask the policy-optimization objective
\citep{gu2025maskdpo}, to restructure the reward model
\citep{qiu2025sentencelevel}, to supply dense rewards during reinforcement
learning \citep{wu2023finegrained}, to train a reward model that outputs
per-step judgments directly \citep{lightman2023verify}, and to build
evaluation metrics \citep{laban2022summac, song2024finesure, min2023factscore}.
What distinguishes this thesis is the
insertion point. Each of the training-oriented works above takes the preference
labels as given, in most cases as response-level human judgments, and modifies
what happens downstream of them. The step examined here is the construction of
those labels: how the output of an automatic factuality evaluator, queried at
different granularities, is converted into a pairwise preference, and how
learnable the resulting labels are when everything downstream is held constant.
This is attractive because automatic labels are cheap, and uncertain because they
are noisier per label than human ones, so whether the additional granularity
recovers more signal than noise is not obvious in advance.

\section{Positioning of This Thesis}
\label{sec:rw-positioning}

The contribution of this thesis is not sentence-level reward modeling itself, nor
sentence-level factuality alignment in general, both of which are active areas
with established results. It is a controlled comparison of
\emph{preference-construction procedures}: the downstream Bradley--Terry
architecture and training procedure are held fixed while the granularity at which
automatic factuality evidence is converted into pairwise preference labels is
varied. No new factuality scorer is introduced; AlignScore is used off the shelf,
and the measurement instrument is the learnability of the resulting labels.

Six adjacent lines of work are distinguished as follows:

\begin{itemize}
  \item \textbf{RLAIF and LLM-as-a-judge research}
    \citep{bai2022constitutional, lee2023rlaif, zheng2023judging, gu2024survey}
    studies judge reliability at the level of a whole-response verdict, rather
    than the effect of querying the judge at a finer granularity.
  \item \textbf{Fine-grained human feedback} \citep{wu2023finegrained}
    establishes that localized supervision helps when labels are reliable; the
    present study asks the cheaper and less certain version, in which each local
    label is produced automatically.
  \item \textbf{Sentence-level factuality alignment} \citep{gu2025maskdpo}
    applies granular factuality information inside the policy-optimization
    objective while leaving preference construction unchanged; here the objective
    is unchanged and preference construction varies.
  \item \textbf{Sentence-level reward modeling} \citep{qiu2025sentencelevel}
    changes the reward model's architecture while learning from conventional
    response-level preferences; here the architecture is fixed and the
    preferences vary.
  \item \textbf{Process reward modeling} \citep{lightman2023verify} trains a
    reward model to output human-labeled, per-step correctness judgments
    directly, consumed downstream for reranking or search; here the reward
    model outputs a single scalar per summary, unchanged from the holistic
    condition, and what varies is only how that scalar's training label was
    constructed upstream.
  \item \textbf{Fine-grained evaluation frameworks}
    \citep{laban2022summac, song2024finesure, min2023factscore} decompose
    summaries in order to diagnose them rather than to supervise with them.
\end{itemize}

The comparison is conducted on preference data generated by one evaluator on one
corpus, so the findings describe behavior under specific conditions rather than a
general property of granular supervision
(Chapter~\ref{ch:discussion}, Section~\ref{sec:disc-generalizability}).
